\section{Review of Related Literature}

\subsection{Automated Feeding and Watering Systems}

Automated feeding and watering systems in commercial piggeries are designed to deliver controlled and consistent amounts of feed and water using sensor-based monitoring and programmable control logic. These systems aim to reduce feed wastage and improve efficiency by replacing manual, labor-dependent practices with precise and repeatable operations \cite{adeoye2025animalia, beyondtech2025smart}.\\

Literature emphasizes that inconsistencies in feeding and drinking behavior can lead to health deterioration and inefficient resource utilization \cite{oecdfao2025outlook, livestockproduction2026}. As a result, automated systems are designed to standardize feeding schedules and quantities while ensuring that consumption is monitored continuously. A key consideration in system design is maintaining consistency across multiple feeding points, which requires coordinated control and reliable data exchange within the system.

\subsection{Health Monitoring Through Behavioral and Consumption Data}

Behavioral monitoring through feeding and drinking patterns is widely recognized as a non-invasive method for early disease detection in swine. Changes in consumption behavior often occur before visible clinical symptoms, making them valuable indicators of potential health issues \cite{dingcong2017monitoring, sadeghi2023improving}.\\

Studies such as Dingcong et al. demonstrate that tracking feeder and drinker interactions can provide insight into animal health and behavior \cite{dingcong2017monitoring}. More broadly, research highlights the importance of establishing baseline consumption patterns and identifying deviations over time \cite{sadeghi2023improving, racewicz2021welfare}. These deviations can be analyzed using time-series approaches, where historical data is used to define normal behavior and detect anomalies.\\

Without structured data collection and analysis, caretakers rely on subjective observation, which may result in delayed detection of illness. The literature supports the use of continuous monitoring systems to provide objective and data-driven assessment of animal health conditions \cite{dingcong2017monitoring, racewicz2021welfare}.

\subsection{Air Quality Sensing in Swine Environments}

Environmental conditions within piggery enclosures play a significant role in animal health and productivity. Poor air quality has been associated with respiratory illnesses and reduced overall welfare in swine populations \cite{livestockproduction2026, yamsakul2025machine}.\\

Research indicates that monitoring environmental parameters enables early identification of unsafe conditions and supports timely intervention. Air quality data also complements behavioral monitoring, as environmental stressors can influence feeding and drinking patterns \cite{sadeghi2023improving, yamsakul2025machine}. This relationship highlights the importance of integrating environmental and behavioral data to achieve a more comprehensive understanding of animal health.\\

Existing studies also explore advanced approaches such as sound-based detection of respiratory symptoms, further reinforcing the role of non-invasive monitoring techniques in livestock management \cite{yamsakul2025machine}.

\subsection{Centralized Data Monitoring and User Interface Design}

Centralized monitoring systems are widely used in agricultural automation to aggregate data from multiple sensing points and provide a unified view of system status. These systems enable real-time monitoring, data logging, and alert generation, supporting informed decision-making by farm operators \cite{beyondtech2025smart, racewicz2021welfare}.\\

The literature emphasizes the importance of accessible and intuitive user interfaces that present complex data in a simplified manner. Features such as real-time dashboards, historical trend visualization, and threshold-based alerts are commonly identified as essential components of effective monitoring systems \cite{sadeghi2023improving}.\\

In addition, centralized systems reduce reliance on manual observation and improve the accuracy and consistency of recorded data. Continuous monitoring technologies have been shown to enhance productivity and reduce labor demands in modern swine farming operations \cite{lo2025rebooting, livestockproduction2026}.

\subsection{African Swine Fever and Technology-Driven Biosecurity}

The impact of African Swine Fever (ASF) on the global and Philippine swine industry has underscored the importance of early detection and preventive monitoring systems. The disease has caused significant livestock losses and economic disruption, prompting increased investment in biosecurity and surveillance initiatives \cite{hsu2025perspective, lo2025rebooting}.\\

Programs such as the Bantay ASF sa Barangay (BABay ASF) and the BRIDGES project focus on monitoring and rapid diagnosis at the community and national levels \cite{hsu2025perspective, pcaarrd2026}. However, literature suggests that farm-level monitoring systems can complement these efforts by providing continuous observation of animal behavior and environmental conditions.\\

By detecting deviations that may indicate early-stage illness, monitoring systems contribute to preventive health management. This approach aligns with broader efforts to reduce reliance on antibiotics and improve sustainability in livestock production \cite{healthforanimals2024prevention}.

\subsection{Scalability and Deployment Considerations in Philippine Piggeries}

The structure of the Philippine swine industry, characterized by a combination of commercial farms and smallholder operations, presents unique challenges for technology adoption. With a large proportion of hog production managed by backyard farmers, systems must be adaptable, scalable, and accessible \cite{lo2025rebooting}.\\

Literature highlights the need for flexible system designs that can accommodate varying farm sizes and operational conditions. Continuous monitoring technologies are increasingly being adopted in modern farming environments, demonstrating their potential to improve productivity and animal welfare \cite{lo2025rebooting, livestockproduction2026}.\\

Additionally, non-invasive monitoring approaches, including sound-based and image-based techniques, provide scalable methods for observing animal health across large populations \cite{reza2025review, dingcong2017monitoring}. These approaches expand the range of available tools for precision livestock farming and support the gradual modernization of the sector.